% Partie : théorie de la démonstration
% Chapitre : calcul des séquents
% Section : traduction des preuves

\documentclass[../../../include/open-logic-section]{subfiles}

\begin{document}

\olfileid[fr]{pt}{seq}{tr}

\olsection{Traduction entre \Log{G1c} et \Log{G3c}}

Nous avons indiqué que \Log{G1c} et \Log{G3c} sont tous deux
corrects et complets pour la logique classique : chacun prouve
exactement les séquents vrais dans toute structure.
Ils prouvent donc les mêmes séquents. Nous pouvons maintenant
établir ce résultat par la seule théorie de la démonstration,
sans passer par les théorèmes de correction et de complétude.
Cette preuve fournit des \emph{transformations} des preuves
de \Log{G1c} en preuves de \Log{G3c}, et réciproquement.

\begin{prop}
Si $\Log{G1c} \Proves S$, alors $\Log{G3c} \Proves S$.
\end{prop}

\begin{proof}
Soit $\pi$ une preuve du séquent~$S$ dans \Log{G1c}.
Nous construisons une preuve~$\pi'$ de~$S$ dans \Log{G3c},
par récurrence sur $n = \pheight{\pi}$.

Si $n = 0$, alors $\pi$ est un axiome.
L'axiome $\lfalse \Sequent \quad$ est aussi un axiome de
\Log{G3c}, en prenant les contextes $\Gamma, \Delta$ vides.
C'est dans \Log{G1c} que les axiomes $!A \Sequent !A$
sont permis pour une formule~$!A$ quelconque.
Dans \Log{G3c}, les axiomes d'identité sont restreints
aux formules atomiques. Cependant, nous avons montré dans la
\olref[ex]{prop:G3c-axioms} que tous ces séquents d'identité
sont prouvables dans~\Log{G3c}.

Si $n>0$, la preuve~$\pi$ se termine par une règle~$R$.
L'hypothèse de récurrence fournit des preuves des prémisses
de cette inférence dans~\Log{G3c} : les sous-preuves immédiates
possèdent donc des traductions. Lorsque la règle~$R$ est
identique dans les deux calculs, on l'applique à ces preuves.

Pour $\LeftR{\land}$ et $\RightR{\lor}$, la
\olref[adm]{prop:lor-G3-adm} établit la dérivabilité des versions
de \Log{G3c} dans \Log{G1c}. Dans le sens requis ici,
on ajoute à la prémisse traduite, par affaiblissement admissible,
le conjoint manquant à gauche, ou le disjoint manquant à droite ;
puis on applique la règle correspondante de~\Log{G3c}.
Pour une règle d'affaiblissement, on applique la
\olref[adm]{prop:weak-G3c-adm}. Pour une règle de contraction,
on applique la \olref[inv]{prop:G3c-cont-adm}.

Il faut aussi traiter les deux règles quantifiées qui diffèrent.
Pour \LeftR{\lforall}, on ajoute à l'antécédent de la prémisse
traduite une occurrence de l'universelle principale, puis
on applique la règle gauche de \Log{G3c}, qui conserve cette
universelle dans sa prémisse. Pour \RightR{\lexists},
on ajoute de même au conséquent une occurrence de
l'existentielle principale, puis on applique la règle droite
de~\Log{G3c}. Ces ajouts utilisent l'affaiblissement admissible,
avec renommage préalable des constantes propres si nécessaire.
Les instances restent celles des mêmes termes clos.
Les deux autres règles quantifiées sont identiques dans les
deux calculs, et leur constante propre reste absente de la
conclusion inchangée. Tous les cas sont ainsi couverts.
\end{proof}

\begin{prop}
Si $\Log{G3c} \Proves S$, alors $\Log{G1c} \Proves S$.
\end{prop}

\begin{proof}
Nous raisonnons de nouveau par récurrence sur la hauteur
d'une preuve~$\pi$ du séquent~$S$.
Chaque axiome de \Log{G3c} se prouve dans~\Log{G1c}
à partir d'un axiome suivi de plusieurs applications
de \LeftR{\Weakening} et de \RightR{\Weakening} :
\[
  \Axiom$!C \fCenter !C$
  \RightLabel{\LeftR{\Weakening}}
  \Deduce$!C, \Gamma \fCenter !C$
  \RightLabel{\RightR{\Weakening}}
  \Deduce$!C, \Gamma \fCenter \Delta, !C$
  \DisplayProof
\qquad
  \Axiom$\lfalse \fCenter$
  \RightLabel{\LeftR{\Weakening}}
  \Deduce$\lfalse, \Gamma \fCenter$
  \RightLabel{\RightR{\Weakening}}
  \Deduce$\lfalse, \Gamma \fCenter \Delta$
  \DisplayProof
  \]
Si $\pi$ se termine par une règle~$R$, l'hypothèse de récurrence
fournit des preuves de ses prémisses dans~\Log{G1c}.
On applique la même règle~$R$ lorsque ses schémas sont identiques.
Les règles \LeftR{\land} et \RightR{\lor} diffèrent entre
\Log{G1c} et~\Log{G3c}. Leurs versions dans~\Log{G3c}
sont dérivables dans~\Log{G1c} par les schémas suivants :
\[
  \Axiom$!A, !B, \Gamma \fCenter \Delta$
  \RightLabel{\LeftR{\land}}
  \UnaryInf$!A \land !B, !B, \Gamma \fCenter \Delta$
  \RightLabel{\LeftR{\land}}
  \UnaryInf$!A \land !B, !A \land !B, \Gamma \fCenter \Delta$
  \RightLabel{\LeftR{\Contraction}}
  \UnaryInf$!A \land !B, \Gamma \fCenter \Delta$
    \DisplayProof
\qquad
  \Axiom$\Gamma \fCenter \Delta, !A, !B$
  \RightLabel{\RightR{\lor}}
  \UnaryInf$\Gamma \fCenter \Delta, !A \lor !B, !B$
  \RightLabel{\RightR{\lor}}
  \UnaryInf$\Gamma \fCenter \Delta, !A \lor !B, !A \lor !B$
  \RightLabel{\RightR{\Contraction}}
  \UnaryInf$\Gamma \fCenter \Delta, !A \lor !B$
    \DisplayProof
  \]
Il reste les deux règles quantifiées qui conservent leur
principale dans la prémisse de \Log{G3c}.
Pour \LeftR{\lforall}, la prémisse traduite contient l'instance
et une occurrence de l'universelle principale. On applique
\LeftR{\lforall} de \Log{G1c} à l'instance en laissant
l'universelle dans le contexte. La conclusion contient alors
deux occurrences de cette universelle ; une contraction gauche
redonne la conclusion recherchée. Pour \RightR{\lexists},
on applique la règle de \Log{G1c} à l'instance en conservant
l'existentielle principale dans le contexte du conséquent.
Les deux occurrences obtenues sont ensuite contractées à droite.
Les termes d'instanciation ne changent pas, et les autres
règles quantifiées se traduisent directement.
\end{proof}

Ces transformations établissent l'existence d'une preuve
dans l'autre calcul. Les étapes ajoutées peuvent augmenter
sa hauteur : aucune conservation de hauteur n'est affirmée.

\end{document}
